\subsection{Nonunique minimizers}\label{sec:exact-nonunique}

Theorem~\ref{thm:transfer} reduces exact output to reaching accuracy
$\varepsilon_S$. The algorithm of Section~\ref{sec:growth} does not reach it
at a rate controlled by $(p,\kappa)$ when the minimizer is not unique, even
with two isolated minimizers and set growth: its grids are graded around one
center, and coarse cells near a distant minimizer keep the corrected minimum
low.

\begin{proposition}[One center cannot certify two distant minimizers]
\label{prop:twocenters}
For $M>0$ let
\[
F_M(x,z)=x^2-2xz+Mz\qquad\text{on }X=[0,M]^2\ (\text{both coordinates
continuous}).
\]
Then $S=\{(0,0),(M,M)\}$, $\OPT=0$, $L_x=2$, $L_z=0$, and
$F_M(v)\ge\frac1{20}\dist(v,S)^2$ on $X$, so the set-growth condition number
is at most $40$; one bag $\{x,z\}$ covers both scopes ($p=2$).
Let $\theta\in(0,1/4]$, $h\in(0,M]$ and let $c\in X$ be any center. For the
graded grids~\eqref{eq:step} on $X$ with center $c$, mesh $h$, grading
$\theta$, and any corrections $d_x=L_xw_x^2/8$ with $L_x\ge2$ and $d_z\ge0$,
\begin{equation}\label{eq:twocenters}
\beta=\min_GQ\le-\max\Bigl\{\frac{\theta^2M^2}{379},\ \frac{h^2}{20}\Bigr\}.
\end{equation}
Consequently, in FG and CT every stage has certified gap
$U-\beta\ge\max\{\theta^2M^2/379,\,h^2/20\}$; a stage certifying a
gap $\varepsilon$ uses grids with at least $M/(48\sqrt\varepsilon)$ nodes in
each coordinate; and any exact-output procedure based on these certificates
that accepts only when $U-\beta<1$, such as EX, builds a table with at least
$M^2/2304$ entries. For integer $M$ the input length is $I=O(\log M)$, so
this work is $2^{\Omega(I)}$.
\end{proposition}

\begin{proof}
\emph{The instance.} $F_M=x^2+z(M-2x)$. For $x<M/2$ the minimum over $z$ is
at $z=0$ with value $x^2$; for $x>M/2$ it is at $z=M$ with value $(M-x)^2$;
at $x=M/2$ the value is $M^2/4$. So $S$ and $\OPT$ are as stated, and the
Hessian $\bigl(\begin{smallmatrix}2&-2\\-2&0\end{smallmatrix}\bigr)$ gives
$L_x=2$, $L_z=0$. The map $(x,z)\mapsto(M-x,M-z)$ preserves $F_M$ (a direct
expansion) and swaps the minimizers, so for growth we may assume
$x\le M/2$. If $x\le M/4$, then $M-2x\ge M/2\ge z/2$ and
$F_M\ge x^2+z^2/2\ge\frac12\norm{(x,z)}^2$. If $M/4<x\le M/2$, then
$F_M\ge x^2\ge M^2/16$ while $\norm{(x,z)}^2\le5M^2/4$. In both cases
$F_M\ge\frac1{20}\dist((x,z),S)^2$.

\emph{The bound~\eqref{eq:twocenters}.} By symmetry assume $M-c_x\ge M/2$
(otherwise use the reflected argument with $z=0$ and the endpoint $0$). Let
$c_x=v_0<v_1<\dots<v_K=M$ be the $x$-nodes on the right of the center,
$D_k=v_k-c_x$, $a=v_{K-1}-v_{K-2}$ (if $K\ge2$) and $c'=M-v_{K-1}$. All steps
except the last are unclipped: $D_{k+1}=(1+\theta)D_k+h$ for $k<K-1$, and
$D_K\le(1+\theta)D_{K-1}+h$. The node $z=M$ is in the $z$-grid, and
$F_M(v,M)=(M-v)^2$. Every correction is nonnegative and
$d_x(v)\ge w_x(v)^2/4$.

If $K=1$, the node $M$ has the adjacent interval $[c_x,M]$ of length at least
$M/2$, so $\beta\le Q(M,M)\le-M^2/16$, and
$M^2/16\ge\max\{\theta^2M^2/379,h^2/20\}$ since $\theta\le1/4$ and $h\le M$.

If $K\ge2$, then $Q(M,M)\le-c'^2/4$ and
$Q(v_{K-1},M)\le c'^2-a^2/4$. If $c'^2\ge a^2/5$ the first is at most
$-a^2/20$, otherwise the second is below $-a^2/20$; so $\beta\le-a^2/20$.
Now $a=h+\theta D_{K-2}\ge h$, and the recursion gives
$D_{K-2}\ge(D_K-(2+\theta)h)/(1+\theta)^2$. If $h\le M/16$, then with
$D_K\ge M/2$ and $\theta\le1/4$,
$D_{K-2}\ge(\frac M2-\frac{9M}{64})\cdot\frac{16}{25}=\frac{23M}{100}$, so
$a\ge\frac{23}{100}\theta M$. If $h>M/16$, then $a\ge h>M/16\ge\theta M/4$.
Hence $a\ge\max\{h,\frac{23}{100}\theta M\}$ and
$\beta\le-\max\{h^2,(0.23\,\theta M)^2\}/20$, which
implies~\eqref{eq:twocenters} because $0.23^2/20>1/379$.

\emph{Consequences.} Filtering keeps both minimizers
(Proposition~\ref{prop:filter}), so every box of FG is $[0,M]^2$, and every
stage is of the form above with $U\ge\OPT=0$; this gives the gap bound. A
stage with $U-\beta\le\varepsilon$ has $\theta\le\sqrt{379\varepsilon}/M$
and $h\le\sqrt{20\varepsilon}$, so every step is at most
$h+\theta M<24\sqrt\varepsilon$, and the side of length at least $M/2$ in each
coordinate needs at least $M/(48\sqrt\varepsilon)$ nodes. The table of the
bag $\{x,z\}$ has at least $M^2/(2304\varepsilon)$ entries. CT stops only
after some stage certifies its target, and an acceptance with
$U-\beta<1$ needs a stage with gap below one.
\end{proof}

The instance is trivial for other methods; the point is that the failure
comes from the single center, not from conditioning. Replacing hulls by
unions of retained cells, on uniform meshes, removes it at the price of a
$\sqrt n$ factor per coordinate.

\paragraph{Algorithm UC$(\varepsilon)$ (uniform cells).}
For each coordinate keep a finite set $\mathcal C_i$ of \emph{retained cells},
closed intervals of positive length; initially $\mathcal C_i=\{[\ell_i,u_i]\}$.
Let $J$ be the least $j\ge0$ with $\frac12nLs^24^{-j}\le\varepsilon$, start
with $\hat x=\ell$, $U=F(\ell)$, and for $j=0,1,\dots,J$ with $h_j=s2^{-j}$:
\begin{enumerate}[label=(\roman*)]
\item \emph{Refine.} Let $\Lambda_{ij}=\{\ell_i+kh_j:k\in\Z_{\ge0}\}$ for
$i\in I_C$ and $\Lambda_{ij}=\{\ell_i+\lceil kh_j\rceil:k\in\Z_{\ge0}\}$ for
$i\in I_Z$. A retained cell $[a,b']$ is split at the points of
$\Lambda_{ij}\cap(a,b')$; the resulting \emph{stage cells} have consecutive
nodes as endpoints, and $G_i$ is the set of all their endpoints. Let
$w_i(v)$ be the largest length of a stage cell with endpoint $v$, ignoring
integer cells of length one, and define $d_i$, $\Delta$, $Q$, $\beta_j$ and $m_i$
as in~\eqref{eq:correction}--\eqref{eq:minmarginal} on $G=\prod_iG_i$.
\item \emph{Solve.} Compute $\beta_j$, a minimizer $y_j$ and all $m_i$ by
Lemma~\ref{lem:dp}; if $F(y_j)<U$, set $\hat x=y_j$, $U=F(y_j)$.
\item \emph{Filter.} Let $\mathcal C_i$ be the set of stage cells $[a,b']$
with $\min\{m_i(a),m_i(b')\}\le U$. Do not take hulls.
\end{enumerate}
Return $\hat x$, $U$ and $\beta_J$. The current domain at stage $j$ is
$X_j=X\cap\prod_i\bigcup\mathcal C_i$.

Since $\Lambda_{i,j-1}\subseteq\Lambda_{ij}$ (for integers,
$\lceil kh_{j-1}\rceil=\lceil 2kh_j\rceil$), every stage cell is either a
segment between consecutive points of $\Lambda_{ij}\cap[\ell_i,u_i]$ or such
a segment clipped at $u_i$. Hence a continuous stage cell has length at most
$h_j$; an integer stage cell has length at most $\lceil h_j\rceil$, and if
that length is at least two then $h_j>1$ and $\lceil h_j\rceil<2h_j$.
Consecutive points of $\Lambda_{ij}$ are at least $h_j$ apart for $i\in I_C$
and at least $\max\{1,\lfloor h_j\rfloor\}\ge h_j/2$ apart for $i\in I_Z$.
Each stage cell contains at most one point of $\Lambda_{i,j+1}$ in its
interior.

\begin{lemma}[Cells: interpolation and safe filtering]\label{lem:cells}
Assume only Definition~\ref{def:curvature}. At every stage, for every
$x\in X_j$ there is a random $Y\in G$ with $\E Q(Y)\le F(x)$ and
$Y_i\in\{a,b'\}$ whenever $x_i$ lies in the stage cell $[a,b']$. Hence
$\beta_j\le\min_{X_j}F$, and $\min\{m_i(a),m_i(b')\}\le F(x)$ for every
$x\in X_j$ with $x_i\in[a,b']$. Every global minimizer lies in every $X_j$,
and $\beta_J\le\OPT\le U$.
\end{lemma}

\begin{proof}
For each $i$ choose a stage cell $[a_i,b_i]$ containing $x_i$. If $x_i$ is
an endpoint put $Y_i=x_i$; otherwise let $Y_i=b_i$ with probability
$(x_i-a_i)/(b_i-a_i)$ and $Y_i=a_i$ otherwise, independently. If $i\in I_Z$
and $b_i-a_i=1$, then $x_i$ is an endpoint. Let
$Z^{k}=(Y_1,\dots,Y_k,x_{k+1},\dots,x_n)\in\bar X$. Conditionally on
$Y_1,\dots,Y_{k-1}$, the function $\psi(t)=F(Z^{k-1}+(t-x_k)e_k)-\frac{L_k}2t^2$
is concave on $[a_k,b_k]$, so Jensen's inequality gives
$\E F(Z^k)\le F(Z^{k-1})+\frac{L_k}2\Var(Y_k)$. Summing over $k$,
$\E F(Y)\le F(x)+\sum_k\frac{L_k}2\Var(Y_k)$. If $Y_k$ is rounded, then
$\Var(Y_k)=(x_k-a_k)(b_k-x_k)\le(b_k-a_k)^2/4$ while both possible values have
$d_k\ge L_k(b_k-a_k)^2/8$; if not, $\Var(Y_k)=0\le d_k$. So
$\E D(Y)\ge\sum_k\frac{L_k}2\Var(Y_k)$ and $\E Q(Y)\le F(x)$. Since
$Q(Y)\ge\beta_j$ and, when $x_i\in[a,b']$ is the chosen cell, $Q(Y)\ge
\min\{m_i(a),m_i(b')\}$, the two bounds follow. A point $x\in X_j$ removed
at stage $j$ has some coordinate whose stage cells containing $x_i$ were all
discarded, so $F(x)>U_j\ge U$; in particular minimizers are never removed,
and by induction $X_J$ contains $S$ and $\beta_J\le\OPT$.
\end{proof}

\begin{theorem}[Uniform cells]\label{thm:cells}
Assume Definition~\ref{def:curvature}. For every $\varepsilon>0$,
UC$(\varepsilon)$ returns a feasible $\hat x$ and $\beta\le\OPT\le
F(\hat x)\le\beta+\varepsilon$ with a certificate (all stage tables and
filtering decisions) whose validity uses only Definition~\ref{def:curvature}.
Suppose moreover that $S$ is finite, $r=\max_i\abs{\pi_i(S)}$, and
$F(x)-\OPT\ge g_S\dist(x,S)^2$ on $X$; put $\kappa_S=\max\{1,L/g_S\}$. Then
every grid of every stage has at most
\begin{equation}\label{eq:cellcount}
K_S=12\,r\,\bigl(2\sqrt{n\kappa_S}+1\bigr)
\end{equation}
nodes per coordinate. For an explicit rational quadratic, UC$(2^{-q})$ uses
$O\bigl(p(N+\abs{\mathcal A})K_S^p\bigr)\poly(I+q)$ bit operations, and
EX with UC in place of CT returns an exact minimizer in
$O\bigl(p(N+\abs{\mathcal A})K_S^p\bigr)\poly(I+\log\kappa_S)$ bit
operations. None of $g_S$, $S$, $r$ is an input.
\end{theorem}

\begin{proof}
\emph{Gap.} Every stage cell that receives a correction has length at most
$h_j$ (continuous) or below $2h_j$ (integer), so
$D(y)\le\sum_iL_i(2h_j)^2/8\le\frac12nLh_j^2$ for every $y\in G$. Hence
$U-\beta_j\le F(y_j)-\beta_j=D(y_j)\le\frac12nLh_j^2$, which is at most
$\varepsilon$ at $j=J$. Validity is Lemma~\ref{lem:cells}.

\emph{Count.} Fix a stage $j$ and a cell $[a,b']$ retained by step~(iii). One
endpoint $v$ has $m_i(v)\le U$; let $z\in G$ attain $m_i(v)$, so $z_i=v$.
By Lemma~\ref{lem:cells}, $\beta_j\le\OPT$, so
$U\le F(y_j)\le\OPT+\frac12nLh_j^2$ and
\[
F(z)-\OPT\le U+D(z)-\OPT\le nLh_j^2 .
\]
Set growth gives $\dist(z,S)^2\le nLh_j^2/g_S\le n\kappa_Sh_j^2$, so
$\abs{v-\sigma}\le\rho:=\sqrt{n\kappa_S}\,h_j$ for some $\sigma\in\pi_i(S)$.
The endpoints of stage cells are points of $\Lambda_{ij}$ or $u_i$. An
interval of length $2\rho$ contains at most $2\rho/\delta+1$ points of
$\Lambda_{ij}$, where $\delta=h_j$ (continuous) or $h_j/2$ (integer), plus
possibly $u_i$; each endpoint belongs to at most two stage cells. So at most
$2r(4\sqrt{n\kappa_S}+2)$ cells are retained, and since each contains at most
one new interior node, the stage-$(j+1)$ grid has at most
$3\cdot2r(4\sqrt{n\kappa_S}+2)=K_S$ nodes. Stage $0$ has two nodes per
coordinate.

\emph{Work and bit length.} By Lemma~\ref{lem:dp} each stage takes
$O(p(N+\abs{\mathcal A})K_S^p)$ table operations, and
$J=O(\log(nLs^2/\varepsilon))=O(I+q)$. Continuous nodes have the form
$\ell_i+ks2^{-j}$ and lie in $[\ell_i,u_i]$, integer nodes are integers, so
all nodes have denominators dividing $\Delta 2^j$ and bit length $O(I+j)$; factor
values, corrections and messages have bit length polynomial in $I+j$. For
exact output, Theorem~\ref{thm:transfer}(a) is reached at
$q\le\poly(I)+\log_2\kappa_S$, as in the proof of Theorem~\ref{thm:exact}.
\end{proof}

\begin{remark}[What grading buys]\label{rem:grading}
With $r=1$, Theorem~\ref{thm:cells} gives accuracy-independent grids of
size $O(\sqrt{n\kappa})$ without trial gradings or caps. The graded
single-center grids of Section~\ref{sec:growth} reduce $\sqrt n$ to
$\log n$, which is what makes Theorem~\ref{thm:approx} fixed-parameter
tractable in $(p,\kappa)$; Proposition~\ref{prop:twocenters} shows that
they pay for this with a dependence on a unique minimizer. Whether exact
output admits an $f(p,\kappa_S,r)\poly(I)$ bound when $S$ is finite, or any
accuracy-independent grid bound when $S$ is a continuum, is open.
\end{remark}

\begin{remark}[Relation to rational-witness results]\label{rem:np}
Corollary~\ref{cor:height} is the box case, with an explicit constant, of a
classical fact: a quadratic program with rational data that attains its
minimum has a global minimizer of polynomial encoding length, obtained from
a rational linear system \cite{Vavasis1990}; Vavasis used it to show that the
decision version of quadratic programming is in NP, and Del Pia, Dey and
Molinaro extended NP membership to mixed-integer quadratic programs with
unbounded integer variables \cite{DelPiaDeyMolinaro2017}. On a bounded
mixed box the integer coordinates have polynomial length, so that extension
is not needed. These results certify $\OPT\le t$. Deciding $\OPT\le t$ for
box QP is NP-complete, and NP-hard already at treewidth two
\cite{DelPiaKhajavirad2026}, so polynomial-size, polynomial-time checkable
certificates of lower bounds $\OPT\ge t$ for all instances would place an
NP-complete problem in coNP. The
certificates of Theorem~\ref{thm:exact} have size $f(p,\kappa)\poly(I)$:
they exist for bounded $(p,\kappa)$ and are checked without trusting
growth. The explicit constants $R$ and $\Omega$ are what make the stopping test
computable.
\end{remark}

